% *** ก่อนส่ง: ลบข้อความตัวอย่าง (\ldots) ทั้งหมด แล้วเขียนเนื้อหาของตนเอง ***
% =============================================================
% บทที่ 3 - วิธีการดำเนินงาน
% แก้ไขเนื้อหาในบทนี้ได้ที่ไฟล์ chapter3.tex
% =============================================================

% อธิบายขั้นตอนและวิธีการดำเนินงาน
% รวมถึงเครื่องมือที่ใช้ในการทำโครงงาน

% -----------------------------------------------------------
\section{ขั้นตอนการดำเนินงาน}
% -----------------------------------------------------------

โครงงานนี้ดำเนินงานตามขั้นตอน ดังนี้

\subsection{ศึกษาปัญหาและรวบรวมความต้องการ}

ศึกษาปัญหาจาก\ldots{} และรวบรวมความต้องการจาก\ldots{}

\subsection{วิเคราะห์และออกแบบระบบ}

วิเคราะห์ความต้องการ และออกแบบระบบโดยใช้\ldots{}

\subsection{พัฒนาระบบ}

พัฒนาระบบตามที่ออกแบบไว้ โดยใช้ภาษา\ldots{} ร่วมกับเฟรมเวิร์ค\ldots{}

\subsection{ทดสอบและประเมินผล}

ทดสอบระบบด้วยวิธี\ldots{} และประเมินผลโดย\ldots{}

\subsection{ปรับปรุงและสรุปผล}

ปรับปรุงระบบตามผลการทดสอบ และจัดทำรายงานสรุป

% -----------------------------------------------------------
\section{เครื่องมือที่ใช้ในการดำเนินงาน}
% -----------------------------------------------------------

\subsection{ฮาร์ดแวร์}

เครื่องมือด้านฮาร์ดแวร์ที่ใช้ในการพัฒนาโครงงานแสดงดังตารางที่ \ref{tab:hardware}

\begin{table}[H]
\centering
\caption{รายการฮาร์ดแวร์ที่ใช้ในการพัฒนา}
\label{tab:hardware}
\begin{tabular}{|c|l|l|}
\hline
\textbf{ลำดับ} & \textbf{รายการ} & \textbf{รายละเอียด} \\
\hline
1 & คอมพิวเตอร์ส่วนบุคคล & CPU\ldots{}, RAM\ldots{} GB, SSD\ldots{} GB \\ \hline
2 & จอภาพ & ขนาด\ldots{} นิ้ว \\ \hline
3 & \ldots & \ldots \\
\hline
\end{tabular}
\end{table}

\subsection{ซอฟต์แวร์}

เครื่องมือด้านซอฟต์แวร์ที่ใช้ในการพัฒนาโครงงานแสดงดังตารางที่ \ref{tab:software}

\begin{table}[H]
\centering
\caption{รายการซอฟต์แวร์ที่ใช้ในการพัฒนา}
\label{tab:software}
\begin{tabular}{|c|l|l|l|}
\hline
\textbf{ลำดับ} & \textbf{ซอฟต์แวร์} & \textbf{เวอร์ชัน} & \textbf{วัตถุประสงค์} \\
\hline
1 & ระบบปฏิบัติการ\ldots & \ldots & สภาพแวดล้อมการพัฒนา \\ \hline
2 & ภาษาโปรแกรม\ldots & \ldots & พัฒนาระบบ \\ \hline
3 & เฟรมเวิร์ค\ldots & \ldots & พัฒนา Frontend/Backend \\ \hline
4 & ระบบจัดการฐานข้อมูล\ldots & \ldots & จัดเก็บข้อมูล \\ \hline
5 & IDE\ldots & \ldots & เครื่องมือเขียนโปรแกรม \\ \hline
6 & Git & \ldots & จัดการเวอร์ชัน \\
\hline
\end{tabular}
\end{table}

% -----------------------------------------------------------
\section{แผนการดำเนินงาน}
% -----------------------------------------------------------

แผนการดำเนินงานตลอดระยะเวลาโครงงานแสดงดังตารางที่ \ref{tab:gantt}

\begin{table}[H]
\centering
\caption{แผนการดำเนินงาน (Gantt Chart)}
\label{tab:gantt}
\begin{tabular}{|c|l|c|c|c|c|c|c|}
\hline
\textbf{ลำดับ} & \textbf{กิจกรรม} & \textbf{เดือน 1} & \textbf{เดือน 2} & \textbf{เดือน 3} & \textbf{เดือน 4} & \textbf{เดือน 5} & \textbf{เดือน 6} \\ \hline
1 & ศึกษาปัญหา & $\bullet$ & & & & & \\ \hline
2 & วิเคราะห์/ออกแบบ & & $\bullet$ & $\bullet$ & & & \\ \hline
3 & พัฒนาระบบ & & & $\bullet$ & $\bullet$ & $\bullet$ & \\ \hline
4 & ทดสอบระบบ & & & & & $\bullet$ & $\bullet$ \\ \hline
5 & สรุปผล/รายงาน & & & & & & $\bullet$ \\ \hline
\end{tabular}
\end{table}
